\documentclass[tikz,border=0pt]{standalone}
\usepackage[T1]{fontenc}
\usepackage{mathptmx}
\usepackage{graphicx}
\usetikzlibrary{arrows.meta}
\definecolor{forest}{HTML}{346E57}
\begin{document}
\begin{tikzpicture}[x=1cm,y=1cm]
  \path[use as bounding box] (0,0) rectangle (16,5);
  \fill[white] (0,0) rectangle (16,5);
  \node[inner sep=0pt] at (2.45,2.9) {\includegraphics[width=4.8cm,height=3.7cm,keepaspectratio]{overview_city.png}};
  \node[inner sep=0pt] at (8,2.9) {\includegraphics[width=4.8cm,height=3.7cm,keepaspectratio]{overview_facility.png}};
  \node[inner sep=0pt] at (13.55,2.9) {\includegraphics[width=4.8cm,height=3.7cm,keepaspectratio]{overview_device.png}};
  \draw[-{Latex[length=2mm,width=1.4mm]},forest,line width=.7pt] (4.98,2.9) -- (5.47,2.9);
  \draw[-{Latex[length=2mm,width=1.4mm]},forest,line width=.7pt] (10.53,2.9) -- (11.02,2.9);
  \node[font=\fontsize{10.5}{12}\selectfont\bfseries,text=black,inner sep=0pt,text width=4.95cm,align=center] at (2.45,.53) {Cocokkan perjalanan};
  \node[font=\fontsize{10.5}{12}\selectfont\bfseries,text=black,inner sep=0pt,text width=4.95cm,align=center] at (8,.53) {Terima dan periksa};
  \node[font=\fontsize{10.5}{12}\selectfont\bfseries,text=black,inner sep=0pt,text width=4.95cm,align=center] at (13.55,.53) {Tentukan jalur pemulihan};
\end{tikzpicture}
\end{document}
